\begin{table}[htbp]\centering\small
\caption{Robustness of the ln GACI coefficient: ln SO2/GDP}\label{tab:A1a}
\begin{tabular}{lccccccccccc}
\toprule
 & base & +K/L & +sea MA & +LSCI 06- & GACI-region$\times$yr & continent$\times$yr & subregion$\times$yr & region trends & tercile$\times$yr & 1996-2009 & 2010-2019 \\
\midrule
ln GACI & -1.024*** & -0.982*** & -0.834** & -0.574* & -0.348 & -0.444 & -0.258 & -0.326 & -0.309 & -0.862*** & -0.371 \\
 & (0.334) & (0.338) & (0.337) & (0.296) & (0.303) & (0.356) & (0.349) & (0.287) & (0.318) & (0.334) & (0.253) \\
\midrule
Observations & 3,818 & 3,551 & 3,818 & 1,788 & 3,753 & 3,818 & 3,813 & 3,753 & 3,520 & 2,170 & 1,647 \\
\bottomrule
\end{tabular}
\begin{minipage}{0.95\textwidth}\footnotesize Unified sample: 173 countries, 1996--2019, country-years with all Table 1 outcomes and the instrument observed. All denominators from one source (WDI GDP in constant 2015 US\$ = GDP per capita $\times$ population), so log identities hold exactly. Country and year fixed effects; controls ln population, ln GDP per capita and its square. Standard errors clustered by country. *, **, *** : 10, 5, 1\%. K/L from PWT 11. GACI-region = the 7 macro-regions of the GACI panel (AF, AS, EU, LA, ME, NA, SW); continent = 5 UN continents; subregion = 22 UN subregions; region trends = region-specific linear trends. CO2 columns 12--13 extend the sample to 2023 (CEDS ends 2019).\end{minipage}
\end{table}